\documentclass[11pt]{article}
\usepackage[utf8]{inputenc}
\usepackage[T1]{fontenc}
\usepackage{amsmath,amssymb}
\usepackage{graphicx}
\usepackage{booktabs}
\usepackage{hyperref}
\usepackage[margin=1in]{geometry}
\usepackage{natbib}
\usepackage{authblk}

\title{Time-, Tissue- and Treatment-Associated Heterogeneity in Tumour-Residing Migratory Dendritic Cells}

\author[1]{Author A}
\author[1]{Author B}
\author[1,2]{Author C}
\affil[1]{Department of Immunology, Weizmann Institute of Science, Rehovot, Israel}
\affil[2]{Department of Oncology, Tel Aviv Sourasky Medical Center, Tel Aviv, Israel}

\date{}

\begin{document}
\maketitle

\begin{abstract}
CCR7$^+$ mature regulatory dendritic cells (mRegDCs) are key orchestrators of anti-tumour immunity, migrating from tumours to draining lymph nodes (dLNs) to activate T cells. However, a subset of mRegDCs remains tumour-resident despite expressing CCR7, and the phenotypic consequences of prolonged tumour dwell-time are unknown. Using photoconvertible mouse models to time-stamp dendritic cells in solid tumours, we tracked migration kinetics and characterized mRegDC heterogeneity by tissue location and dwell-time. We find that mRegDCs constitute 78--85\% of DCs migrating to the dLN, yet 18\% of photoconverted DCs remain in the tumour at 72 hours. Tumour-retained mRegDCs exhibit a regulatory phenotype with 2-fold higher PD-L1 expression, 0.59-fold lower MHC-II, and 0.41-fold lower costimulatory molecules (CD80, IL-12) compared to dLN counterparts. Single-cell RNA sequencing reveals progressive downregulation of antigen presentation and pro-inflammatory gene modules with increasing dwell-time, while regulatory programs intensify. Tumour mRegDCs co-localize with PD-1$^+$CD8$^+$ T cells at 3.2--3.6-fold above random expectation in both murine and human tumours. Anti-PD-L1 treatment partially reverses the regulatory phenotype, increasing MHC-II expression by 1.5-fold and IL-12$^+$ mRegDC frequency from 12\% to 28\%, while boosting CD8$^+$ T cell infiltration from 180 to 420 cells/mm$^2$. These findings reveal that the tumour microenvironment progressively reprograms resident mRegDCs toward an immunosuppressive state that can be partially reversed by checkpoint immunotherapy.
\end{abstract}

\section{Introduction}

Dendritic cells (DCs) are professional antigen-presenting cells that bridge innate and adaptive immunity. In the tumour microenvironment, a specialized population of CCR7$^+$ mature DCs---termed mature regulatory DCs (mRegDCs) or mregDCs---has been identified as a critical mediator of anti-tumour immune responses \citep{maier2020ctimp, gerhard2021tumor}. These cells acquire tumour antigens, upregulate CCR7 and costimulatory molecules, and migrate to tumour-draining lymph nodes (dLNs) where they activate tumour-specific T cells.

However, CCR7 expression does not guarantee emigration from the tumour. A subset of mRegDCs remains tumour-resident despite expressing the migration receptor, and this population is poorly understood \citep{maier2020ctimp}. The tumour microenvironment is rich in immunosuppressive signals---including PD-L1, IL-10, TGF-$\beta$, and prostaglandins---that could alter the phenotype and function of DCs that fail to emigrate in a timely fashion. Whether prolonged tumour dwell-time progressively reprograms mRegDCs and how this affects their interactions with tumour-infiltrating T cells are open questions with important implications for immunotherapy.

Anti-PD-L1/PD-1 checkpoint immunotherapy has transformed cancer treatment, but its effects on DC biology within the tumour microenvironment are not fully characterized. PD-L1 is highly expressed by mRegDCs \citep{maier2020ctimp}, suggesting that these cells could be direct targets of anti-PD-L1 therapy.

Here, we use photoconvertible mouse models to precisely time-stamp DCs within tumours and track their migration kinetics. By combining flow cytometry, single-cell RNA sequencing, and spatial analysis, we characterize the heterogeneity of mRegDCs by tissue location and dwell-time. We find that tumour-retained mRegDCs undergo progressive regulatory reprogramming and that anti-PD-L1 treatment partially reverses this phenotype.

\section{Related Work}

The identification of mRegDCs as a distinct DC population in tumours was established by \citet{maier2020ctimp}, who used single-cell RNA sequencing to define a CCR7$^+$ Fscn1$^+$ population expressing both immunostimulatory and regulatory molecules. \citet{gerhard2021tumor} further characterized the functional significance of tumour-associated DCs and their role in shaping anti-tumour immunity.

Photoconvertible mouse models (Kaede and KikGR) have been used to study immune cell migration in lymph nodes and skin \citep{tomura2008monitoring}, but their application to tracking DC dwell-time dynamics within tumours has been limited. The temporal evolution of DC phenotype as a function of tissue residence time has not been systematically characterized.

The spatial relationship between DCs and T cells in the tumour microenvironment has been explored using multiplexed imaging approaches \citep{schurch2020coordinated}. Co-localization of antigen-presenting cells with CD8$^+$ T cells has been associated with favorable immunotherapy outcomes, but the specific role of mRegDCs in this spatial organization has not been established.

The effects of checkpoint immunotherapy on tumour DC biology have been studied primarily at the population level \citep{garris2018successful}. Whether anti-PD-L1 treatment specifically alters the regulatory phenotype of tumour-resident mRegDCs has not been directly examined.

\section{Methods}

\subsection{Photoconvertible mouse models and tumour implantation}
Kaede transgenic mice, expressing a photoconvertible fluorescent protein that switches from green to red upon UV illumination, were used for DC migration tracking. MC38 colon carcinoma or B16 melanoma cells were implanted subcutaneously. At day 10 post-implantation, tumours were exposed to UV light to photoconvert all cells in situ. Converted (red) DCs were tracked in the tumour and dLN at 6, 24, 48, and 72 hours post-conversion.

\subsection{Flow cytometry}
Multi-parameter flow cytometry panels were designed to identify mRegDCs (CD11c$^+$ MHC-II$^+$ CCR7$^+$) and assess expression of functional markers including PD-L1, CD80, CD86, and intracellular IL-12. Median fluorescence intensity (MFI) values were compared between tumour and dLN mRegDCs.

\subsection{Single-cell RNA sequencing}
Tumour mRegDCs were sorted by photoconversion status and time point, then profiled by 10x Genomics Chromium scRNA-seq. Cells were binned by dwell-time: short ($<$24 h), medium (24--48 h), and long ($>$48 h). Differential expression was calculated as log$_2$ fold change relative to freshly arrived cells.

\subsection{Spatial analysis}
Multiplexed immunohistochemistry was performed on murine (MC38, B16) and human (melanoma, NSCLC) tumour sections to identify mRegDCs and PD-1$^+$CD8$^+$ T cells. Co-localization was defined as mRegDCs within 30 $\mu$m of PD-1$^+$CD8$^+$ T cells and compared to random expectation.

\subsection{Anti-PD-L1 treatment}
MC38 tumour-bearing mice were treated with anti-PD-L1 antibody or isotype control starting at day 10. Tumour mRegDCs and CD8$^+$ T cell infiltration were profiled at day 5 post-treatment initiation.

\section{Results}

\subsection{mRegDCs dominate tumour-to-dLN migration but a subset remains tumour-resident}
Photoconversion tracking revealed that mRegDCs constitute the dominant population migrating from the tumour to the dLN, representing 78--85\% of converted DCs arriving in the dLN across all time points (Table~\ref{tab:migration}). However, a substantial fraction of photoconverted DCs remained in the tumour even at 72 hours post-conversion (18\% of converted DCs), indicating that a subset of mRegDCs fails to emigrate despite CCR7 expression.

\begin{table}[h]
\centering
\caption{DC migration kinetics post-photoconversion.}
\label{tab:migration}
\begin{tabular}{cccc}
\toprule
Time post-conversion & \% converted DCs in tumour & \% converted DCs in dLN & \% mRegDC among dLN arrivals \\
\midrule
6 h & 85 & 5 & 78 \\
24 h & 52 & 28 & 82 \\
48 h & 30 & 42 & 85 \\
72 h & 18 & 38 & 80 \\
\bottomrule
\end{tabular}
\end{table}

\subsection{Tumour-retained mRegDCs have a regulatory phenotype}
Flow cytometric comparison of tumour and dLN mRegDCs at 48 hours post-photoconversion revealed striking phenotypic differences (Table~\ref{tab:phenotype}). Tumour mRegDCs expressed 2.0-fold higher PD-L1 ($p < 0.001$) but 0.59-fold lower MHC-II, 0.41-fold lower CD80, 0.51-fold lower CD86, and 0.41-fold lower intracellular IL-12 compared to dLN counterparts. This indicates that tumour-retained mRegDCs have reduced antigen-presenting and pro-inflammatory capacity but elevated regulatory molecule expression.

\begin{table}[h]
\centering
\caption{Phenotypic comparison of tumour versus dLN mRegDCs at 48 h post-photoconversion.}
\label{tab:phenotype}
\begin{tabular}{lcccc}
\toprule
Marker (MFI) & Tumour mRegDC & dLN mRegDC & Fold difference & $p$-value \\
\midrule
CCR7 & 2,800 & 4,500 & 0.62x & $<$0.01 \\
PD-L1 & 6,200 & 3,100 & 2.0x & $<$0.001 \\
MHC-II & 3,400 & 5,800 & 0.59x & $<$0.001 \\
CD80 & 1,200 & 2,900 & 0.41x & $<$0.001 \\
CD86 & 1,800 & 3,500 & 0.51x & $<$0.01 \\
IL-12 (intracellular) & 450 & 1,100 & 0.41x & $<$0.01 \\
\bottomrule
\end{tabular}
\end{table}

\subsection{Progressive regulatory reprogramming with increased dwell-time}
scRNA-seq analysis of tumour mRegDCs binned by dwell-time revealed progressive transcriptomic changes (Table~\ref{tab:temporal}). Antigen presentation genes (H2-Aa, H2-Ab1) showed progressive downregulation ($-0.8$ log$_2$FC at medium dwell, $-1.6$ at long dwell). Pro-inflammatory genes (Il12b, Tnf) declined even more steeply ($-1.1$ and $-2.3$ log$_2$FC). Conversely, regulatory genes (Cd274/PD-L1, Socs2) were progressively upregulated ($+0.5$ and $+1.2$ log$_2$FC). Migration-associated genes (Ccr7, Fscn1) also declined with dwell-time, consistent with the observed failure to emigrate.

\begin{table}[h]
\centering
\caption{Expression changes in tumour mRegDCs by dwell-time (log$_2$FC vs.\ freshly arrived).}
\label{tab:temporal}
\begin{tabular}{lccc}
\toprule
Gene module & Short ($<$24 h) & Medium (24--48 h) & Long ($>$48 h) \\
\midrule
Antigen presentation & 0 (ref) & $-$0.8 & $-$1.6 \\
Pro-inflammatory & 0 (ref) & $-$1.1 & $-$2.3 \\
Regulatory & 0 (ref) & $+$0.5 & $+$1.2 \\
Migration & 0 (ref) & $-$0.3 & $-$0.9 \\
\bottomrule
\end{tabular}
\end{table}

\subsection{Spatial co-localization with PD-1$^+$CD8$^+$ T cells}
Multiplexed immunohistochemistry revealed significant spatial co-localization between tumour mRegDCs and PD-1$^+$CD8$^+$ T cells (Table~\ref{tab:spatial}). Across both murine (MC38, B16) and human (melanoma, NSCLC) tumour types, 32--41\% of mRegDCs were within 30 $\mu$m of PD-1$^+$CD8$^+$ T cells, representing a 3.2--3.6-fold enrichment over random expectation. This spatial proximity suggests functional interactions between regulatory mRegDCs and exhausted T cells within the tumour microenvironment.

\begin{table}[h]
\centering
\caption{Spatial co-localization of mRegDCs with PD-1$^+$CD8$^+$ T cells.}
\label{tab:spatial}
\begin{tabular}{lccc}
\toprule
Tumour model & mRegDC proximity (\%) & Random expectation (\%) & Enrichment fold \\
\midrule
Mouse (MC38) & 38 & 11 & 3.5x \\
Mouse (B16) & 32 & 9 & 3.6x \\
Human (melanoma) & 41 & 13 & 3.2x \\
Human (NSCLC) & 35 & 10 & 3.5x \\
\bottomrule
\end{tabular}
\end{table}

\subsection{Anti-PD-L1 treatment partially reverses the regulatory phenotype}
Treatment of MC38 tumour-bearing mice with anti-PD-L1 antibody significantly altered tumour mRegDC biology (Table~\ref{tab:treatment}). Tumour mRegDC frequency increased from 4.2\% to 6.8\% of CD45$^+$ cells ($p < 0.01$). MHC-II expression increased by approximately 1.5-fold, and the proportion of IL-12$^+$ mRegDCs more than doubled (12\% to 28\%; $p < 0.001$). PD-L1 expression on mRegDCs was reduced from MFI 6,100 to 2,800 ($p < 0.001$). Concurrently, CD8$^+$ T cell infiltration increased from 180 to 420 cells/mm$^2$ ($p < 0.001$).

\begin{table}[h]
\centering
\caption{Effects of anti-PD-L1 treatment on tumour mRegDCs and CD8$^+$ T cell infiltration.}
\label{tab:treatment}
\begin{tabular}{lccc}
\toprule
Parameter & Isotype control & Anti-PD-L1 & $p$-value \\
\midrule
mRegDC frequency (\% CD45$^+$) & 4.2 & 6.8 & $<$0.01 \\
MHC-II MFI & 3,200 & 4,900 & $<$0.01 \\
IL-12$^+$ mRegDC (\%) & 12 & 28 & $<$0.001 \\
PD-L1 MFI & 6,100 & 2,800 & $<$0.001 \\
CD8$^+$ T cells (cells/mm$^2$) & 180 & 420 & $<$0.001 \\
\bottomrule
\end{tabular}
\end{table}

\section{Discussion}

Our study reveals that the tumour microenvironment imposes a time-dependent regulatory reprogramming on resident mRegDCs, converting them from immunostimulatory antigen-presenting cells into PD-L1$^\text{hi}$ regulatory cells with diminished capacity to activate T cells. This process is progressive: the longer mRegDCs remain in the tumour, the more pronounced the regulatory phenotype becomes.

The spatial co-localization of tumour mRegDCs with PD-1$^+$CD8$^+$ T cells, observed in both murine and human tumours, suggests that these regulatory mRegDCs may directly suppress T cell function through the PD-L1/PD-1 axis. The 3.2--3.6-fold enrichment over random expectation indicates an active spatial organization, possibly driven by chemokine gradients or antigen-dependent interactions.

The finding that anti-PD-L1 treatment partially reverses the regulatory phenotype of tumour mRegDCs has important therapeutic implications. Beyond blocking the PD-L1/PD-1 inhibitory signal between mRegDCs and T cells, anti-PD-L1 therapy appears to remodel mRegDC biology itself, restoring MHC-II expression and IL-12 production. This suggests that the therapeutic benefit of checkpoint blockade may extend beyond direct T cell reinvigoration to include restoration of DC function \citep{garris2018successful}.

The photoconversion approach provides temporal resolution not achievable with static single-cell analyses, enabling us to track phenotypic changes as a function of residence time. However, a limitation is that photoconversion is irreversible and labels all cells in the tumour, not just DCs; DC identity must be established post hoc by flow cytometry or sorting. Additionally, the overexpression of the photoconvertible protein itself could theoretically affect DC biology, although we observed no evidence of this.

\section{Conclusion}

Using photoconvertible mouse models, we demonstrate that mRegDCs are the dominant DC population migrating from tumours to draining lymph nodes, but a subset remains tumour-resident and undergoes progressive regulatory reprogramming characterized by upregulation of PD-L1 and downregulation of antigen presentation and pro-inflammatory capacity. These tumour-resident mRegDCs co-localize with PD-1$^+$CD8$^+$ T cells, and anti-PD-L1 treatment partially reverses their regulatory phenotype while boosting T cell infiltration. These findings identify tumour-resident mRegDCs as targets and mediators of checkpoint immunotherapy.

\bibliographystyle{plainnat}
\bibliography{references}

\end{document}
